\chapter{§2.1 有理数}
\label{sec:2.1}


\prerequisite{§1.1 整数的认识, §1.5 分数, §1.6 小数}

\gradelevel{7 年级}


\section*{负数的引入}


\begin{explore}


冬天的哈尔滨，天气预报播报明日最低气温为"零下 15 度"。在地理课上，老师说死海的海拔是"低于海平面 430 米"。银行账户里存入 500 元后又取出 800 元，余额变成了"欠款 300 元"。

这些场景中，"零下""低于""欠款"都表达了一个方向——与通常的正方向\textbf{相反}。仅靠 0 和正数，我们无法精确记录这些量。数学需要一种新的数。

\end{explore}

\begin{understand}


\textbf{直觉理解}：在日常生活中，很多量都有两个相反的方向：盈利与亏损、上升与下降、零上与零下。我们约定一个方向为"正"，相反方向就是"负"。

\textbf{定义}：

\begin{itemize}
  \item \textbf{正数}：大于 0 的数，如 $1, 2.5, \frac{3}{4}$ 。
  \item \textbf{负数}：在正数前加负号" $-$ "，如 $-1, -2.5, -\frac{3}{4}$ 。
  \item $0$ 既不是正数，也不是负数。
\end{itemize}


\textbf{为什么需要负数？} 如果只有正数和零，我们无法表示"低于零度""亏损"等现实中的量。引入负数后，每一个量都可以用一个数精确描述，数学的表达能力大大增强。

\end{understand}

\begin{workedexamples}


\textbf{例 1.} 用正数或负数表示下列各量：

(1) 零上 $8^\circ\text{C}$ 　　(2) 零下 $3^\circ\text{C}$ 　　(3) 海拔 $+200$ 米　　(4) 海拔 $-50$ 米

\textbf{解：}

(1) 零上 $8^\circ\text{C}$ 记为 $+8^\circ\text{C}$ （或 $8^\circ\text{C}$ ）。

(2) 零下 $3^\circ\text{C}$ 记为 $-3^\circ\text{C}$ 。

(3) 海拔 $+200$ 米表示高于海平面 $200$ 米。

(4) 海拔 $-50$ 米表示低于海平面 $50$ 米。

\textbf{例 2.} 某公司一月份盈利 $20$ 万元，二月份亏损 $5$ 万元，三月份盈利 $12$ 万元。用正负数记录每月的经营情况。

\textbf{解：} 约定盈利为正、亏损为负。

\begin{itemize}
  \item 一月份： $+20$ （万元）
  \item 二月份： $-5$ （万元）
  \item 三月份： $+12$ （万元）
\end{itemize}


\textbf{例 3.} 小明乘电梯，从地面（0 层）上升 $5$ 层，再下降 $8$ 层到达地下停车场。用正负数表示两次位移，并求最终位置。

\textbf{解：} 约定上升为正、下降为负。

\begin{itemize}
  \item 第一次位移： $+5$ （层）
  \item 第二次位移： $-8$ （层）
  \item 最终位置： $5 + (-8) = -3$ （层），即地下第 $3$ 层。
\end{itemize}


\end{workedexamples}

\begin{keytakeaway}


\begin{itemize}
  \item 正数和负数用来表示意义相反的量。
  \item $0$ 是正数和负数的分界，既不是正数也不是负数。
  \item 引入负数后，数轴上 $0$ 的左边有了无穷多个数。
\end{itemize}


\end{keytakeaway}

\begin{practice}


\begin{enumerate}
  \item 用正数或负数表示：收入 $1500$ 元；支出 $800$ 元。
  \item 如果向东走 $3$ 千米记为 $+3$ 千米，那么向西走 $7$ 千米记为多少？
  \item 海平面的海拔记为 $0$ 米。某地海拔 $-155$ 米，这表示什么？
  \item 某天北京最高气温为 $5^\circ\text{C}$ ，最低气温为 $-7^\circ\text{C}$ ，这一天的温差是多少？
  \item 判断对错： $0$ 是最小的正数。
\end{enumerate}


\end{practice}

\section*{有理数的分类}


\begin{explore}


我们已经认识了正整数（ $1, 2, 3, \ldots$ ）、零、负整数（ $-1, -2, -3, \ldots$ ）以及正分数（ $\frac{1}{2}, 0.3, \ldots$ ）。引入负数之后，分数也有了负数形式，如 $-\frac{1}{2}, -0.3$ 。这么多类型的数，有没有统一的分类方式？

\end{explore}

\begin{understand}


\textbf{有理数}是整数和分数的统称。按\textbf{符号}分类：

\[
\text{有理数}\begin{cases}
\text{正有理数}\begin{cases} \text{正整数：} 1, 2, 3, \ldots \\ \text{正分数：} \frac{1}{2}, 0.7, \ldots \end{cases} \\[6pt]
\text{零：} 0 \\[6pt]
\text{负有理数}\begin{cases} \text{负整数：} -1, -2, -3, \ldots \\ \text{负分数：} -\frac{1}{2}, -0.7, \ldots \end{cases}
\end{cases}
\]


也可以按\textbf{是否为整数}分类：

\[
\text{有理数}\begin{cases}
\text{整数}\begin{cases} \text{正整数} \\ 0 \\ \text{负整数} \end{cases} \\[6pt]
\text{分数}\begin{cases} \text{正分数} \\ \text{负分数} \end{cases}
\end{cases}
\]


\textbf{为什么叫"有理数"？} 英文 rational number 来自 ratio（比），有理数就是能写成两个整数之比 $\frac{p}{q}$ （ $q \neq 0$ ）的数。

\end{understand}

\begin{workedexamples}


\textbf{例 1.} 把下列各数填入相应的集合中：

\[
-7,\; \frac{2}{3},\; 0,\; -3.5,\; 100,\; -\frac{1}{4},\; 0.8,\; -11
\]


正整数集合： $\{\ \}$ ；负整数集合： $\{\ \}$ ；正分数集合： $\{\ \}$ ；负分数集合： $\{\ \}$ 。

\textbf{解：}

\begin{itemize}
  \item 正整数集合： $\{100\}$
  \item 负整数集合： $\{-7, -11\}$
  \item 正分数集合： $\{\frac{2}{3}, 0.8\}$
  \item 负分数集合： $\{-3.5, -\frac{1}{4}\}$
\end{itemize}


注意： $0$ 是整数，但既不是正整数也不是负整数。

\textbf{例 2.} 判断： $-3.0$ 是整数还是分数？

\textbf{解：} $-3.0 = -3$ ，所以 $-3.0$ 是负整数。

\textbf{例 3.} $\frac{6}{2}$ 是分数还是整数？

\textbf{解：} $\frac{6}{2} = 3$ ，化简后是正整数。虽然写成了分数的形式，但实质上是整数。

\end{workedexamples}

\begin{keytakeaway}


\begin{itemize}
  \item 有理数 = 整数 + 分数。
  \item 分类时注意 $0$ 的特殊地位：它是整数，但不属于正数或负数。
  \item 判断一个数的类别，先化简再分类。
\end{itemize}


\end{keytakeaway}

\begin{practice}


\begin{enumerate}
  \item $-\frac{5}{1}$ 是整数还是分数？
  \item 把 $3, -\frac{7}{2}, 0, 1.5, -9, \frac{4}{2}$ 按正整数、负整数、正分数、负分数分类。
  \item 判断对错：所有分数都是有理数。
  \item 判断对错： $0$ 是最小的整数。
  \item 有理数中，最大的负整数是多少？最小的正整数是多少？
\end{enumerate}


\end{practice}

\section*{数轴、相反数、绝对值}


\begin{explore}


在温度计上， $0$ 度的刻度线把温度计分成了上下两部分：上面是正温度，下面是负温度。每个温度对应温度计上的唯一一个点。这种"用一条线上的点来表示数"的想法，正是\textbf{数轴}的来源。

\end{explore}

\begin{understand}


\textbf{数轴}：规定了\textbf{原点}（表示 $0$ ）、\textbf{正方向}（通常向右）和\textbf{单位长度}的直线，叫做数轴。

数轴上的每一个点都对应一个有理数，每一个有理数都对应数轴上的一个点。

\textbf{相反数}：如果两个数只有符号不同，我们就说它们互为\textbf{相反数}。例如 $3$ 和 $-3$ 互为相反数。 $0$ 的相反数是 $0$ 本身。

\begin{itemize}
  \item 数 $a$ 的相反数是 $-a$ 。
  \item $a + (-a) = 0$ 。
\end{itemize}


\textbf{绝对值}：数轴上表示一个数的点到原点的\textbf{距离}，叫做这个数的\textbf{绝对值}，记为 $\lvert a \rvert$ 。

\[
\lvert a \rvert = \begin{cases} a, & \text{若 } a > 0 \\ 0, & \text{若 } a = 0 \\ -a, & \text{若 } a < 0 \end{cases}
\]


\textbf{为什么绝对值总是非负的？} 因为距离不可能是负数。 $\lvert -5 \rvert = 5$ 表示 $-5$ 离原点的距离是 $5$ 个单位。

\end{understand}

\begin{workedexamples}


\textbf{例 1.} 在数轴上标出下列各数： $-3, -1.5, 0, 2, 4.5$ 。

\textbf{解：} 画一条数轴，取适当的单位长度：

\[
\underset{-3}{\bullet} \quad \underset{-1.5}{\bullet} \quad \underset{0}{\bullet} \quad \underset{2}{\bullet} \quad \underset{4.5}{\bullet}
\]


（在数轴上， $-3$ 在原点左侧 $3$ 个单位处， $-1.5$ 在原点左侧 $1.5$ 个单位处， $2$ 在原点右侧 $2$ 个单位处， $4.5$ 在原点右侧 $4.5$ 个单位处。）

\textbf{例 2.} 求下列各数的相反数和绝对值： $-7$ ， $\frac{3}{4}$ ， $0$ 。

\textbf{解：}

\begin{center}
\begin{tabular}{lll}
\toprule
数 & 相反数 & 绝对值 \\
\midrule
$-7$ & $7$ & $\lvert -7 \rvert = 7$ \\
$\frac{3}{4}$ & $-\frac{3}{4}$ & $\lvert \frac{3}{4} \rvert = \frac{3}{4}$ \\
$0$ & $0$ & $\lvert 0 \rvert = 0$ \\
\bottomrule
\end{tabular}
\end{center}


\textbf{例 3.} 已知 $\lvert x \rvert = 5$ ，求 $x$ 。

\textbf{解：} 绝对值为 $5$ 的数离原点距离为 $5$ ，在数轴上有两个点： $x = 5$ 或 $x = -5$ 。

\textbf{例 4.} 比较大小： $-3$ 与 $-8$ 。

\textbf{解：} 在数轴上， $-3$ 在 $-8$ 的右边，所以 $-3 > -8$ 。

也可以这样理解：两个负数比大小，绝对值小的反而大。 $\lvert -3 \rvert = 3 < 8 = \lvert -8 \rvert$ ，所以 $-3 > -8$ 。

\end{workedexamples}

\begin{keytakeaway}


\begin{itemize}
  \item 数轴是数与几何的桥梁：右边的数大于左边的数。
  \item 相反数： $a$ 与 $-a$ 互为相反数，它们在数轴上关于原点对称。
  \item 绝对值是"到原点的距离"，永远非负。
  \item 比较两个负数大小：绝对值小的反而大。
\end{itemize}


\end{keytakeaway}

\begin{practice}


\begin{enumerate}
  \item 写出 $-12$ 、 $\frac{5}{6}$ 、 $-2.7$ 的相反数。
  \item 求 $\lvert -9 \rvert$ 、 $\lvert 3.14 \rvert$ 、 $\lvert 0 \rvert$ 。
  \item 若 $\lvert a \rvert = 3$ ， $a$ 可以是哪些数？
  \item 比较大小： $-4$ 与 $-1$ ； $-\frac{1}{2}$ 与 $-\frac{3}{4}$ 。
  \item 数轴上 $A$ 点表示 $-2$ ， $B$ 点表示 $5$ ， $A$ 、 $B$ 两点之间的距离是多少？
  \item 若 $\lvert a \rvert = a$ ，则 $a$ 满足什么条件？
\end{enumerate}


\end{practice}

\section*{有理数的四则运算}


\begin{explore}


小华在做一道计算题：一月份的气温变化为先下降 $3^\circ\text{C}$ ，再上升 $7^\circ\text{C}$ ，再下降 $2^\circ\text{C}$ 。如果初始温度为 $5^\circ\text{C}$ ，最终温度是多少？这需要计算 $5 + (-3) + 7 + (-2)$ 。包含负数的加减法该怎么算？

\end{explore}

\begin{understand}


\textbf{有理数的加法法则}：

\begin{enumerate}
  \item \textbf{同号相加}：取相同的符号，绝对值相加。
\end{enumerate}

   - 例： $(-3) + (-5) = -(3+5) = -8$
\begin{enumerate}
  \item \textbf{异号相加}：取绝对值较大的那个数的符号，用较大绝对值减去较小绝对值。
\end{enumerate}

   - 例： $(-7) + 4 = -(7-4) = -3$
   - 例： $(-3) + 8 = +(8-3) = 5$
\begin{enumerate}
  \item \textbf{一个数加零}：等于它本身。 $a + 0 = a$ 。
\end{enumerate}


\textbf{有理数的减法法则}：减去一个数等于加上这个数的相反数。

\[
a - b = a + (-b)
\]


\textbf{有理数的乘法法则}：

\begin{enumerate}
  \item \textbf{同号得正，异号得负}，绝对值相乘。
\end{enumerate}

   - $(-3) \times (-4) = +12$
   - $(-3) \times 4 = -12$
\begin{enumerate}
  \item \textbf{任何数乘 $0$ 等于 $0$ }。
\end{enumerate}


\textbf{为什么负负得正？} 可以从"方向的方向"理解： $(-1) \times a$ 表示将 $a$ 取反方向。那么 $(-1) \times (-1)$ 就是将"反方向"再取反，回到了正方向。

\textbf{有理数的除法法则}：除以一个不为零的数等于乘以它的倒数。

\[
a \div b = a \times \frac{1}{b} \quad (b \neq 0)
\]


符号规则与乘法相同：同号得正，异号得负。

\end{understand}

\begin{workedexamples}


\textbf{例 1.} 计算 $(-6) + 10$ 。

\textbf{解：} 异号相加，绝对值大的是 $10$ （正数）， $10 - 6 = 4$ ，取正号。

\[
(-6) + 10 = 4
\]


\textbf{例 2.} 计算 $(-5) - (-3)$ 。

\textbf{解：} 减去 $-3$ 等于加上 $3$ 。

\[
(-5) - (-3) = (-5) + 3 = -(5 - 3) = -2
\]


\textbf{例 3.} 计算 $(-4) \times (-7)$ 。

\textbf{解：} 同号得正， $4 \times 7 = 28$ 。

\[
(-4) \times (-7) = 28
\]


\textbf{例 4.} 计算 $(-24) \div 6$ 。

\textbf{解：} 异号得负， $24 \div 6 = 4$ 。

\[
(-24) \div 6 = -4
\]


\end{workedexamples}

\begin{keytakeaway}


\begin{itemize}
  \item 加法：同号取公共符号加绝对值；异号取大绝对值的符号减绝对值。
  \item 减法转化为加法： $a - b = a + (-b)$ 。
  \item 乘除法的符号规则：同号得正，异号得负。
  \item "负负得正"是乘法中最核心也最容易出错的规则。
\end{itemize}


\end{keytakeaway}

\begin{practice}


\begin{enumerate}
  \item 计算 $(-8) + 5$ 。
  \item 计算 $(-3) + (-9)$ 。
  \item 计算 $7 - (-4)$ 。
  \item 计算 $(-6) \times 5$ 。
  \item 计算 $(-2) \times (-3) \times (-4)$ 。
  \item 计算 $(-36) \div (-9)$ 。
  \item 计算 $0 \div (-5)$ 。
  \item 如果温度从 $-3^\circ\text{C}$ 上升了 $10^\circ\text{C}$ ，现在的温度是多少？
\end{enumerate}


\end{practice}

\section*{有理数的混合运算}


\begin{explore}


期末考试后，老师以班级平均分 $75$ 分为基准，记录了五位同学的成绩差值： $+12, -8, +5, -3, 0$ 。如果要计算这五位同学的平均分与班级平均分之间的差，需要计算 $(12 + (-8) + 5 + (-3) + 0) \div 5$ 。当加减乘除混合在一起，运算顺序怎么确定？

\end{explore}

\begin{understand}


有理数混合运算的\textbf{运算顺序}与正数中学过的相同：

\begin{enumerate}
  \item 先算\textbf{乘方}（后续将学习），再算\textbf{乘除}，最后算\textbf{加减}。
  \item 有括号的先算括号内，括号从\textbf{小到大}依次计算：小括号 → 中括号 → 大括号。
  \item 同级运算从\textbf{左到右}依次进行。
\end{enumerate}


\textbf{乘法分配律在负数中同样成立}：

\[
a \times (b + c) = a \times b + a \times c
\]


灵活运用运算律可以简化计算。

\end{understand}

\begin{workedexamples}


\textbf{例 1.} 计算 $(-3) \times 4 + (-18) \div 6$ 。

\textbf{解：} 先算乘除，再算加减。

\[
\begin{aligned}
(-3) \times 4 + (-18) \div 6 &= (-12) + (-3) \\
&= -15
\end{aligned}
\]


\textbf{例 2.} 计算 $(-2) \times [3 - (-4)] + 5$ 。

\textbf{解：} 先算中括号内，再算乘法，最后算加法。

\[
\begin{aligned}
(-2) \times [3 - (-4)] + 5 &= (-2) \times [3 + 4] + 5 \\
&= (-2) \times 7 + 5 \\
&= -14 + 5 \\
&= -9
\end{aligned}
\]


\textbf{例 3.} 用简便方法计算 $(-25) \times 36 + (-25) \times 64$ 。

\textbf{解：} 利用分配律提取公因数。

\[
\begin{aligned}
(-25) \times 36 + (-25) \times 64 &= (-25) \times (36 + 64) \\
&= (-25) \times 100 \\
&= -2500
\end{aligned}
\]


\textbf{例 4.} 计算 $\frac{5}{6} - \frac{1}{3} + \frac{7}{12} - \frac{3}{4}$ 。

\textbf{解：} 统一为分母 $12$ 的分数再计算。

\[
\begin{aligned}
\frac{5}{6} - \frac{1}{3} + \frac{7}{12} - \frac{3}{4} &= \frac{10}{12} - \frac{4}{12} + \frac{7}{12} - \frac{9}{12} \\
&= \frac{10 - 4 + 7 - 9}{12} \\
&= \frac{4}{12} = \frac{1}{3}
\end{aligned}
\]


\end{workedexamples}

\begin{keytakeaway}


\begin{itemize}
  \item 混合运算的顺序：先乘除后加减，有括号先算括号。
  \item 运算律（交换律、结合律、分配律）对有理数同样成立，善于使用可简化计算。
  \item 每一步都注意符号，尤其在"减去负数"和"负数相乘"时。
\end{itemize}


\end{keytakeaway}

\begin{practice}


\begin{enumerate}
  \item 计算 $(-2) \times 5 + 3 \times (-4)$ 。
  \item 计算 $(-48) \div 8 - (-25) + (-16)$ 。
  \item 计算 $(-3) \times [(-4) + 6] \div (-2)$ 。
  \item 用简便方法计算 $(-8) \times 125 \times (-0.5)$ 。
  \item 计算 $\frac{1}{2} + \left(-\frac{2}{3}\right) - \left(-\frac{1}{6}\right)$ 。
  \item 某公司连续 $4$ 天的利润变化为 $+5$ （万元）、 $-3$ （万元）、 $-2$ （万元）、 $+8$ （万元），求 $4$ 天的总利润变化。
\end{enumerate}


\end{practice}



\begin{answer}


\textbf{负数的引入 练一练}

\begin{enumerate}
  \item 收入 $1500$ 元记为 $+1500$ 元，支出 $800$ 元记为 $-800$ 元。
  \item $-7$ 千米。
  \item 该地海拔低于海平面 $155$ 米。
  \item $5 - (-7) = 5 + 7 = 12^\circ\text{C}$ 。
  \item 错。 $0$ 既不是正数也不是负数，不是"最小的正数"。
\end{enumerate}


\textbf{有理数的分类 练一练}

\begin{enumerate}
  \item $-\frac{5}{1} = -5$ ，是负整数。
  \item 正整数： $\{3, \frac{4}{2}\}$ （ $\frac{4}{2}=2$ ）；负整数： $\{-9\}$ ；正分数： $\{1.5\}$ ；负分数： $\{-\frac{7}{2}\}$ 。 $0$ 是整数，单独归入整数。
  \item 对。所有分数（正分数和负分数）都是有理数。
  \item 错。整数中没有最小的数，负整数可以无限小。
  \item 最大的负整数是 $-1$ ，最小的正整数是 $1$ 。
\end{enumerate}


\textbf{数轴、相反数、绝对值 练一练}

\begin{enumerate}
  \item $-12$ 的相反数是 $12$ ； $\frac{5}{6}$ 的相反数是 $-\frac{5}{6}$ ； $-2.7$ 的相反数是 $2.7$ 。
  \item $\lvert -9 \rvert = 9$ ， $\lvert 3.14 \rvert = 3.14$ ， $\lvert 0 \rvert = 0$ 。
  \item $a = 3$ 或 $a = -3$ 。
  \item $-4 < -1$ ； $-\frac{1}{2} > -\frac{3}{4}$ （因为 $\frac{1}{2} < \frac{3}{4}$ ，绝对值小的负数反而大）。
  \item $\lvert 5 - (-2) \rvert = \lvert 7 \rvert = 7$ ， $A$ 、 $B$ 之间的距离是 $7$ 。
  \item $\lvert a \rvert = a$ 说明 $a \geq 0$ 。
\end{enumerate}


\textbf{有理数的四则运算 练一练}

\begin{enumerate}
  \item $(-8) + 5 = -3$ 。
  \item $(-3) + (-9) = -12$ 。
  \item $7 - (-4) = 7 + 4 = 11$ 。
  \item $(-6) \times 5 = -30$ 。
  \item $(-2) \times (-3) \times (-4) = 6 \times (-4) = -24$ 。（三个负数相乘，结果为负。）
  \item $(-36) \div (-9) = 4$ 。
  \item $0 \div (-5) = 0$ 。
  \item $-3 + 10 = 7^\circ\text{C}$ 。
\end{enumerate}


\textbf{有理数的混合运算 练一练}

\begin{enumerate}
  \item $(-2) \times 5 + 3 \times (-4) = -10 + (-12) = -22$ 。
  \item $(-48) \div 8 - (-25) + (-16) = -6 + 25 + (-16) = 3$ 。
  \item $(-3) \times [(-4) + 6] \div (-2) = (-3) \times 2 \div (-2) = -6 \div (-2) = 3$ 。
  \item $(-8) \times 125 \times (-0.5) = (-8) \times (-0.5) \times 125 = 4 \times 125 = 500$ 。
  \item $\frac{1}{2} + \left(-\frac{2}{3}\right) - \left(-\frac{1}{6}\right) = \frac{1}{2} - \frac{2}{3} + \frac{1}{6} = \frac{3}{6} - \frac{4}{6} + \frac{1}{6} = 0$ 。
  \item $5 + (-3) + (-2) + 8 = 8$ （万元）， $4$ 天总利润增加 $8$ 万元。
\end{enumerate}


\end{answer}
